% Part: methods
% Chapter: induction
% Section: structural-induction

\documentclass[../../../include/open-logic-section]{subfiles}

\begin{document}

\olfileid[id]{mth}{ind}{sti}

\olsection{Induksi Struktural}

Sejauh ini kita telah menggunakan induksi untuk menetapkan hasil mengenai
semua bilangan asli. Namun, suatu prinsip yang bersesuaian dapat digunakan
secara langsung untuk membuktikan hasil mengenai semua !!{element} suatu
himpunan yang didefinisikan secara induktif. Prinsip ini sering disebut
induksi \emph{struktural}, karena bergantung pada struktur objek-objek yang
didefinisikan secara induktif.

Secara umum, suatu definisi induktif diberikan oleh (a)~daftar !!{element}
``awal'' dari himpunan tersebut dan (b)~daftar operasi yang menghasilkan
!!{element} baru dari himpunan itu berdasarkan anggota lama. Dalam kasus
suku rapi, misalnya, objek awalnya adalah huruf-huruf. Kita hanya mempunyai
satu operasi, yakni
\begin{align*}
  o(s_1, s_2) = & [s_1 \circ s_2].
\end{align*}
Bahkan, Anda dapat memandang bilangan asli~$\Nat$ sendiri sebagai sesuatu
yang diberikan oleh definisi induktif: objek awalnya adalah~$0$, dan
operasinya adalah fungsi penerus~$x + 1$.

Untuk membuktikan sesuatu mengenai semua anggota suatu himpunan yang
didefinisikan secara induktif, yakni bahwa setiap !!{element} himpunan itu
memiliki sifat~$P$, kita harus:
\begin{enumerate}
\item membuktikan bahwa objek-objek awal memiliki~$P$;
\item membuktikan bahwa bagi setiap operasi~$o$, jika argumen-argumennya
  memiliki~$P$, hasilnya juga demikian.
\end{enumerate}
Misalnya, untuk membuktikan sesuatu mengenai semua suku rapi, kita akan
membuktikan bahwa hal itu benar bagi semua huruf dan bahwa hal itu benar
bagi $[s_1 \circ s_2]$ asalkan benar bagi $s_1$ dan $s_2$ masing-masing.

\begin{prop}
  Banyaknya $[$ sama dengan banyaknya $]$ dalam sebarang suku rapi~$t$.
\end{prop}

\begin{proof}
Kita menggunakan induksi struktural. Suku rapi didefinisikan secara
induktif, dengan huruf sebagai objek awal dan operasi $o$ untuk membangun
suku rapi baru dari suku rapi lama.
\begin{enumerate}
\item Klaim tersebut benar bagi setiap huruf, sebab banyaknya $[$ dalam
  sebuah huruf dengan sendirinya adalah~$0$, dan banyaknya $]$ di dalamnya
  juga~$0$.
\item Andaikan banyaknya $[$ dalam $s_1$ sama dengan banyaknya $]$, dan hal
  yang sama benar bagi $s_2$. Banyaknya $[$ dalam $o(s_1, s_2)$, yakni
  dalam $[s_1 \circ s_2]$, adalah jumlah banyaknya $[$ dalam $s_1$ dan
  $s_2$, ditambah satu. Banyaknya $]$ dalam $o(s_1, s_2)$ adalah jumlah
  banyaknya $]$ dalam $s_1$ dan $s_2$, ditambah satu. Jadi, banyaknya $[$
  dalam $o(s_1, s_2)$ sama dengan banyaknya $]$ dalam $o(s_1,s_2)$.
\end{enumerate}
\end{proof}

\begin{prob}
  Buktikan melalui induksi struktural bahwa tidak ada suku rapi yang diawali
  dengan~$]$.
\end{prob}

Mari kita berikan bukti lain melalui induksi struktural: suatu segmen awal
sejati dari untai simbol~$t$ adalah sebarang untai~$s$ yang bersesuaian
dengan $t$ simbol demi simbol ketika dibaca dari kiri, tetapi $t$ lebih
panjang. Jadi, misalnya, $[a \circ {}$ merupakan segmen awal sejati dari
$[a \circ b]$, sedangkan $[b \circ {}$ bukan (keduanya berbeda pada simbol
kedua), dan $[a \circ b]$ juga bukan (keduanya sama panjang).

\begin{prop}\ollabel{prop:initial}
  Setiap segmen awal sejati dari suatu suku rapi~$t$ memiliki lebih banyak
  $[$ daripada $]$.
\end{prop}

\begin{proof}
  Melalui induksi pada~$t$:
  \begin{enumerate}
  \item $t$ hanya berupa satu huruf: maka $t$ tidak mempunyai segmen awal
    sejati.
  \item $t = [s_1 \circ s_2]$ untuk suatu suku rapi $s_1$ dan~$s_2$. Jika
    $r$ merupakan segmen awal sejati dari $t$, terdapat sejumlah
    kemungkinan:
    \begin{enumerate}
    \item $r$ hanya berupa $[$: maka $r$ memiliki satu $[$ lebih banyak
      daripada~$]$.
    \item $r$ berupa $[r_1$, dengan $r_1$ merupakan segmen awal sejati
      dari~$s_1$: karena $s_1$ adalah suku rapi, menurut hipotesis induksi
      $r_1$ memiliki lebih banyak $[$ daripada $]$, dan hal yang sama benar
      bagi $[r_1$.
    \item $r$ berupa $[s_1$ atau $[s_1 \circ {}$: menurut hasil sebelumnya,
      banyaknya $[$ dan $]$ dalam~$s_1$ sama; jadi, banyaknya $[$ dalam
      $[s_1$ ataupun $[s_1 \circ {}$ adalah satu lebih banyak daripada
      banyaknya~$]$.
    \item $r$ berupa $[s_1 \circ r_2$, dengan $r_2$ merupakan segmen awal
      sejati dari~$s_2$: menurut hipotesis induksi, $r_2$ memuat lebih
      banyak $[$ daripada $]$. Menurut hasil sebelumnya, banyaknya $[$ dan
      $]$ dalam~$s_1$ sama. Jadi, banyaknya $[$ dalam $[s_1 \circ r_2$
      lebih besar daripada banyaknya~$]$.
    \item $r$ berupa $[s_1 \circ s_2$: menurut hasil sebelumnya, banyaknya
      $[$ dan $]$ dalam $s_1$ sama, dan begitu pula dalam~$s_2$. Jadi,
      banyaknya $[$ dalam $[s_1 \circ s_2$ adalah satu lebih banyak daripada
      banyaknya $]$.
    \end{enumerate}
  \end{enumerate}
\end{proof}

\end{document}
